\chapter{Image-Guided Pain Procedures}

\section*{Patient focus}
This chapter explains the examination or clinical problem in accessible language. Patients should use it to prepare questions and understand the purpose, experience, benefits, and limitations of imaging; local instructions and the official report take priority.


\label{ch:ch82_pain_procedures}

Image-guided pain procedures use fluoroscopy (real-time X-ray) or CT to precisely deliver medications (steroid, local anaesthetic) or perform denervation to relieve chronic pain. They are indicated when conservative management (analgesics, physiotherapy, lifestyle modification) has failed and a specific pain generator has been identified.

---

\section{1. Patient Perspective: What a Pain Injection Involves}

If you have chronic back, neck, or joint pain that has not improved with tablets and physiotherapy, your doctor may suggest a \textbf{targeted pain injection}.

\begin{itemize}
  \item The procedure is performed under local anaesthetic with light sedation if needed.
  \item Using X-ray (fluoroscopy) or CT guidance, the doctor directs a thin needle to the exact source of your pain.
  \item A mixture of local anaesthetic (for immediate relief) and steroid (for long-term anti-inflammatory effect) is injected.
  \item You may feel pressure during the injection but should not feel sharp pain.
  \item You will rest for 30--60 minutes afterwards. Most patients go home the same day.
  \item The steroid may take 2--7 days to become effective. The local anaesthetic provides immediate but temporary relief, which helps confirm the pain source.
  \item You may experience a brief ``flare'' of pain for 1--3 days before improvement begins.
\end{itemize}

---

\section{2. Types of Pain Procedures}

\begin{table}[h]
\centering
\begin{tabular}{l l l}
\hline
\textbf{Procedure} & \textbf{Target} & \textbf{Indication} \\
\hline
Epidural steroid injection & Epidural space & Radiculopathy, spinal stenosis \\
Facet joint injection & Zygapophyseal joint & Facet-mediated neck/back pain \\
Medial branch block & Medial branch nerve & Diagnostic block before RFA \\
Nerve root block & Specific spinal nerve root & Focal radicular pain \\
Sacroiliac joint injection & SI joint & SI joint-mediated low back pain \\
Trigger point injection & Myofascial trigger point & Myofascial pain syndrome \\
\hline
\end{tabular}
\caption{Common Image-Guided Pain Procedures}
\end{table}

\section{3. Fluoroscopic vs CT Guidance}

\begin{table}[h]
\centering
\begin{tabular}{l l l}
\hline
\textbf{Feature} & \textbf{Fluoroscopy} & \textbf{CT Guidance} \\
\hline
Best for & Facet joints, epidural, SI joint & Deep structures, complex anatomy \\
Contrast confirmation & Yes (arthrogram for joints) & Yes (CT myelography) \\
Patient position & Prone or lateral & Prone \\
Radiation & Lower (pulsed) & Higher (serial scans) \\
Real-time & Yes (continuous) & Intermittent (scan-advance-scan) \\
\hline
\end{tabular}
\caption{Fluoroscopy vs CT Guidance for Pain Procedures}
\end{table}

\begin{pearlbox}[Contrast Confirmation]
Always confirm needle position with contrast injection before delivering medication. For facet joint injections, intra-articular contrast confirms intra-articular placement (the ``arthrogram sign''). For epidural injections, contrast confirms epidural spread and excludes intravascular or intrathecal injection.
\end{pearlbox}

---

\section{4. Steroid vs Anaesthetic: What Is Injected}

\begin{itemize}
  \item \textbf{Local anaesthetic} (e.g., bupivacaine 0.25\% or lidocaine 1\%): Provides immediate, temporary relief lasting 1--6 hours. If the patient experiences significant relief during this window, it confirms the targeted structure as the pain generator (diagnostic value).
  \item \textbf{Corticosteroid} (e.g., triamcinolone, dexamethasone, betamethasone): Provides longer-lasting anti-inflammatory effect, typically onset at 2--7 days, peak at 2--4 weeks, duration 4--12 weeks.
  \item \textbf{Combination}: Most procedures use a mixture of both. The anaesthetic provides diagnostic information; the steroid provides therapeutic benefit.
\end{itemize}

\section{5. Facet Joint Injections and Medial Branch Blocks}

\begin{itemize}
  \item \textbf{Facet joint injection}: Intra-articular injection of steroid and anaesthetic into the zygapophyseal joint. Both diagnostic and therapeutic.
  \item \textbf{Medial branch block}: Injection of anaesthetic (with or without steroid) onto the medial branch nerve that innervates the facet joint. Performed as a \textbf{diagnostic} procedure before radiofrequency ablation (RFA).
  \item \textbf{Diagnostic criteria}: $>$ 80\% pain relief after medial branch block confirms the facet joint as the pain generator and indicates the patient will benefit from RFA.
  \item \textbf{Two blocks required}: Two separate medial branch blocks with different anaesthetics (e.g., lidocaine and bupivacaine) are recommended to improve diagnostic specificity before proceeding to RFA.
\end{itemize}

\section{6. Epidural Steroid Injections}

\begin{itemize}
  \item \textbf{Interlaminar approach}: Needle advanced through the ligamentum flavum into the posterior epidural space. Most common approach for cervical and lumbar epidurals.
  \item \textbf{Transforaminal approach}: Needle advanced through the neural foramen, targeting the ventral epidural space near the nerve root. More targeted for focal radiculopathy.
  \item \textbf{Caudal approach}: Needle through the sacral hiatus into the epidural space. Useful for multi-level stenosis or failed prior approaches.
  \item \textbf{Indications}: Radiculopathy, spinal stenosis, failed back surgery syndrome.
  \item \textbf{Limitations}: Maximum 3 injections per year; diminishing returns with repeated injections.
\end{itemize}

\section{7. When to Recommend Repeat vs Different Procedure}

\begin{itemize}
  \item \textbf{Repeat the same procedure} if the first injection provided $>$ 50\% relief but the effect has worn off, and the pain pattern is consistent.
  \item \textbf{Try a different procedure} if the first injection provided $<$ 50\% relief (suggesting the wrong pain generator was targeted), the pain pattern has changed, or new symptoms have developed.
  \item \textbf{Consider radiofrequency ablation} if medial branch blocks are positive --- RFA provides longer-lasting denervation (6--12 months) than steroid injection.
  \item \textbf{Escalate to surgical referral} if pain procedures fail, there is progressive neurological deficit, or imaging shows a surgically correctable lesion.
\end{itemize}

\begin{pearlbox}[Documentation Requirements]
Every pain procedure report must document: (1) indication and clinical history, (2) consent, (3) guidance modality and settings, (4) target level and approach, (5) needle position confirmed with contrast, (6) medications and volumes injected, (7) contrast spread pattern, (8) complications, and (9) post-procedure plan.
\end{pearlbox}

---

\section{8. Structured Reporting Example}

\begin{testbox}[Indications:]
52-year-old female with chronic axial low back pain for 18 months. MRI shows mild degenerative changes at L4-5 and L5-S1 facets. Failed 6 weeks of physiotherapy and NSAIDs. Requesting L4-5 and L5-S1 facet joint injections.
\end{testbox}
\begin{refbox}[Findings:]
\textbf{Procedure}: Under fluoroscopic guidance with the patient prone, the L4-5 and L5-S1 facet joints were identified bilaterally. After local anaesthesia, a 25G spinal needle was advanced into each joint under AP and oblique fluoroscopy. Intra-articular contrast confirmed correct needle position (arthrogram confirmed bilaterally at both levels). A mixture of 1 mL bupivacaine 0.25\% and 1 mL triamcinolone (40 mg) was injected into each joint. Total: 4 injections (bilateral L4-5 and L5-S1).\\
\textbf{Contrast Spread}: Bilateral intra-articular contrast spread confirmed at L4-5 and L5-S1. No vascular or epidural reflux.\\
\textbf{Complications}: None. Patient tolerated the procedure well.
\end{refbox}
\begin{pearlbox}[Impression:]
1. Bilateral L4-5 and L5-S1 facet joint injections performed under fluoroscopic guidance.
2. Intra-articular needle position confirmed with contrast at all four levels.
3. Medications: bupivacaine 0.25\% and triamcinolone 40 mg per joint.
4. Monitor response. If good relief, consider medial branch blocks followed by RFA for longer-term management.
\end{pearlbox}

\clearpage
